% Gerado a partir de docs/tcc.md por scripts/sincronizar_tcc.py.
\begin{resumo}
% Gerado a partir de docs/tcc.md por scripts/sincronizar_tcc.py.
Este trabalho apresenta a concepção e o desenvolvimento inicial de um sistema de avaliação contextual de serviços públicos digitais. O problema investigado é a dificuldade de relacionar avaliações declaradas pelo cidadão ao contexto de navegação em que são solicitadas. O percurso combinou revisão bibliográfica e normativa, Lean Inception e prototipação incremental, organizados pelas atividades da Design Science Research. Foi implementado um protótipo local em Python, com API, páginas simuladas, regras operacionais sobre cliques e duração informada, formulário de feedback e painel administrativo. A telemetria fornece indícios comportamentais, sem medir qualidade percebida. A versão funcional atual registra satisfação por estrelas, atributos declarados e facilidade percebida por um item CES opcional de 1 a 7. Os itens atuais de satisfação e CES foram definidos após o estudo de interface e ainda não foram avaliados com participantes. Separadamente da implementação funcional, oito participantes avaliaram fluxos de interface no Figma, com registro automático de cliques e tempos pelo Useberry e questionário pós-sessão. Quatro participantes realizaram cada fluxo, com alocação auxiliada por inteligência artificial (IA), descrita como sorteio sem registro auditável, e com popup de redirecionamento presente em ambos. Os resultados descrevem conclusão de tarefa, facilidade percebida e relatos pós-sessão. Não comprovam maior participação, maior especificidade dos comentários do formulário, funcionamento dos gatilhos da API ou superioridade do fluxo contextual. A contribuição investigativa está na formulação de uma hipótese comparativa e na distinção entre evidências comportamentais e declaradas; a contribuição tecnológica está na implementação que permite relacioná-las. Foram especificados requisitos, rastreabilidade, critérios de qualidade e protocolos prospectivos para cidadãos, gestores e especialistas. A continuidade exige executar esses protocolos e calibrar as regras antes de uso real.

\vspace{\onelineskip}
\noindent \textbf{Palavras-chave:} serviços públicos digitais; feedback contextual; avaliação de serviços; experiência do usuário; Design Science Research.
\end{resumo}
